\section{Evaluación y Optimización de Modelos de Visión Artificial}

La validación cuantitativa y la preparación para el despliegue en producción
representan la etapa final en el desarrollo de un sistema de percepción
vehicular. En el control de acceso a estacionamientos, los modelos de aprendizaje
profundo no solo deben alcanzar un alto grado de exactitud matemática, sino que
deben operar bajo estrictas restricciones de latencia en dispositivos de cómputo
embebidos. Esta sección examina las métricas estandarizadas para evaluar la
detección, el reconocimiento secuencial y el seguimiento temporal, así como las
técnicas de compresión de modelos necesarias para garantizar su rendimiento en
tiempo real en el borde (\textit{Edge AI}).

\subsection{Métricas de clasificación y detección: Precision, Recall, F1-Score,
mAP}

Como establecen Sokolova y Lapalme \cite{sokolova2009performance} y Padilla et
al. \cite{padilla2021metrics}, la evaluación de clasificadores y detectores se
fundamenta en la matriz de confusión, la cual categoriza las predicciones frente
a las anotaciones de referencia (\textit{ground truth}) en cuatro cuadrantes:
Verdaderos Positivos ($TP$), Falsos Positivos ($FP$), Falsos Negativos ($FN$) y
Verdaderos Negativos ($TN$). 

En la operación real de un parqueadero inteligente, la interpretación de estos
errores adquiere un significado crítico para la seguridad y la calidad del
servicio:
\begin{itemize}
    \item \textbf{Falso Positivo ($FP$):} Ocurre cuando el detector predice la
        presencia de un vehículo o matrícula donde no existe, o cuando clasifica
        erróneamente un automóvil no registrado como autorizado. En la práctica,
        un $FP$ representa una \textbf{falla de seguridad grave}, pues implica
        abrir automáticamente la barrera de acceso a un vehículo no permitido o
        registrar un ingreso inexistente.
        
    \item \textbf{Falso Negativo ($FN$):} Ocurre cuando el modelo omite la
        presencia de un automóvil en el carril o no logra detectar la matrícula
        de un cliente registrado. Operativamente, un $FN$ genera congestión en
        la barrera, ya que deja esperando injustificadamente al usuario
        autorizado y requiere la intervención manual del operador.
\end{itemize}

A partir de estas categorías, se definen las métricas fundamentales de evaluación:

\begin{equation}
    \text{Precision}
    =
    \frac{TP}{TP + FP}, \quad \text{Recall}
    =
    \frac{TP}{TP + FN}, \quad \text{F1-Score}
    =
    2 \cdot \frac{\text{Precision} \cdot \text{Recall}}{\text{Precision} + \text{Recall}}
\end{equation}

Mientras que la \textit{Precision} mide la validez de las aperturas de barrera
efectuadas (evitando falsos accesos), el \textit{Recall} mide la tasa de
recuperación de vehículos autorizados (evitando demoras innecesarias). El
\textit{F1-Score} constituye la media armónica entre ambas, proporcionando un
indicador sintético del desempeño global.

Para evaluar la localización espacial de las cajas delimitadoras
(\textit{bounding boxes}) de los automóviles y sus matrículas, Padilla et al.
\cite{padilla2021metrics} formulan la precisión media promedio 
(\textit{Mean Average Precision}, mAP). El mAP evalúa el área bajo la curva de
Precision-Recall considerando distintos umbrales de Intersección sobre Unión 
(IoU). Típicamente, el $\text{mAP}@0.5$ evalúa la capacidad del detector para
encerrar de forma aproximada el vehículo o la placa (IoU $\ge 0.5$), mientras que
el $\text{mAP}@0.5:0.95$ promedia el desempeño en diez umbrales estrictos,
midiendo la precisión geométrica del ajuste de la caja delimitadora.

\subsection{Métricas de reconocimiento de secuencias: Exactitud por carácter y CER}

En el módulo de Reconocimiento Óptico de Caracteres (OCR), la evaluación
trasciende la clasificación de imágenes individuales para medir la fidelidad en
la lectura de cadenas alfanuméricas continuas. Como detallan Shi et al.
\cite{shi2017crnn} y Auza Banegas et al. \cite{auzabanegas2026blpr}, la métrica
principal para evaluar la precisión del OCR es la Tasa de Error por Carácter
(\textit{Character Error Rate}, CER), la cual se fundamenta en la distancia de
edición de Levenshtein:

\begin{equation}
    \text{CER} = \frac{S + D + I}{N}
\end{equation}

donde $S$ representa el número de sustituciones de caracteres, $D$ las deleciones
(caracteres omitidos), $I$ las inserciones (caracteres fantasma añadidos) y $N$
la cantidad total de caracteres reales en la matrícula de referencia. 

En la lectura de matrículas, el impacto de una sustitución es directo. Por
ejemplo, si el motor de OCR confunde el carácter alfabético \texttt{``O''} con el
dígito numérico \texttt{``0''} en una placa real (prediciendo \texttt{``12O4ABC'}
en lugar de \texttt{``1204ABC''}), la distancia de Levenshtein contabiliza una
sustitución ($S = 1$). Aunque la precisión a nivel de carácter sea elevada
($6/7 = 85.7\%$), la métrica de coincidencia exacta (\textit{exact match})
cae a cero. Dado que las consultas en la base de datos del parqueadero requieren
una coincidencia alfanumérica exacta para autorizar el paso, minimizar el CER
mediante reglas de corrección sintáctica es indispensable para garantizar el
funcionamiento del sistema.

\subsection{Métricas de seguimiento de objetos: MOTA, MOTP e IDF1}

Para evaluar la estabilidad del módulo de seguimiento multiobjeto a lo largo del
tiempo, Bernardin y Stiefelhagen \cite{bernardin2008clear} y Ristani et al.
\cite{ristani2016performance} formalizaron el conjunto de métricas CLEAR MOT:

\begin{enumerate}
    \item \textbf{Precisión de Seguimiento Multiobjeto (\textit{Multiple Object
        Tracking Precision}, MOTP):} Mide la precisión geométrica promedio en la
        estimación de las posiciones de los vehículos coincidentes a lo largo de
        la secuencia de video:
        \begin{equation}
            \text{MOTP} = \frac{\sum_{t, i} d_{t, i}}{\sum_t c_t}
        \end{equation}
        donde $d_{t, i}$ representa la distancia o desalineación espacial entre
        la caja predicha y la real para la coincidencia $i$ en el fotograma $t$,
        y $c_t$ es el número total de coincidencias.

    \item \textbf{Exactitud de Seguimiento Multiobjeto (\textit{Multiple Object
        Tracking Accuracy}, MOTA):} Evalúa la capacidad global del sistema para
        mantener trayectorias continuas, integrando los tres tipos de errores
        posibles:

        \begin{equation}
            \text{MOTA} = 1 - \frac{\sum_t (FN_t + FP_t + IDSW_t)}{\sum_t GT_t}
        \end{equation}
        donde $GT_t$ es la cantidad de vehículos reales en el fotograma $t$, y
        $IDSW_t$ representa el número de cambios de identidad (\textit{Identity
        Switches}). En el parqueadero, un cambio de ID erróneo ($IDSW$) ocurre
        cuando el rastreador asigna un nuevo número de seguimiento a un automóvil
        que avanza lentamente hacia la barrera, provocando registros duplicados
        del mismo evento de ingreso.

    \item \textbf{Puntuación IDF1 (\textit{Identification F1-Score}):} Mide la
        proporción de detecciones correctamente identificadas a lo largo de toda
        la trayectoria completa del vehículo, evaluando la consistencia del mapeo
        de identidades a largo plazo frente a fragmentaciones u oclusiones
        temporales \cite{ristani2016performance}.
\end{enumerate}

\subsection{Compromiso entre exactitud y latencia: Rendimiento en el Borde}

En aplicaciones industriales de control de accesos, el desempeño de un algoritmo
no se juzga únicamente por su exactitud en conjuntos de datos estáticos, sino por
su capacidad para responder bajo límites temporales estrictos en el entorno
operativo. Para garantizar una experiencia fluida al usuario y evitar
embotellamientos en la vía pública, el pipeline completo de visión (detección de
vehículo, localización de placa, lectura de OCR y decisión de barrera) debe
ejecutarse con una \textbf{latencia total inferior a 1 segundo por evento}
\cite{auzabanegas2026blpr}.

El análisis experimental de rendimiento evalúa el compromiso (\textit{trade-off})
entre exactitud y velocidad mediante dos indicadores clave:
\begin{itemize}
    \item \textbf{Tiempo de Inferencia ($\text{ms/img}$):} La latencia individual
        requerida por la red neuronal para procesar una imagen de entrada en la
        GPU o acelerador de cómputo.

    \item \textbf{Fotogramas por Segundo (FPS):} La tasa de refresco continua que
        puede sostener el pipeline en tiempo de ejecución.
\end{itemize}

El desafío reside en que los dispositivos de borde (\textit{Edge AI}) instalados
en las casetas del parqueadero (tales como módulos NVIDIA Jetson u ordenadores
embebidos) disponen de restricciones severas de consumo energético, memoria VRAM
y capacidad de cómputo en comparación con servidores en la nube. Por ello, la
arquitectura del detector debe equilibrar el número de parámetros y la cantidad
de operaciones de punto flotante (GFLOPs) para sostener tasas superiores a 30 FPS
en resolución de alta definición.

\subsection{Compresión de modelos: Cuantización de pesos (FP32, FP16, INT8)}

Para adaptar redes neuronales profundas y complejas a hardware embebido
restringido sin sacrificar su precisión predictiva, Sapkota et al.
\cite{sapkota2026yolo26} y Jocher et al. \cite{jocher2026yolo26} destacan la
aplicación de técnicas de compresión de modelos, siendo la \textbf{cuantización
de pesos} la más efectiva para acelerar la inferencia.

La cuantización consiste en mapear la representación numérica de los parámetros y
activaciones de la red desde formatos continuos de alta precisión hacia formatos
discretos de menor ancho de banda de memoria:
\begin{itemize}
    \item \textbf{FP32 (Punto Flotante de 32 bits):} Es el formato de precisión
        simple estándar utilizado durante el entrenamiento en GPU, ofreciendo un
        amplio rango dinámico a costa de una elevada huella de memoria y mayor
        consumo energético.

    \item \textbf{FP16 (Media Precisión de 16 bits):} Reduce a la mitad el tamaño
        del modelo en memoria RAM/VRAM y acelera la inferencia en aceleradores
        modernos con unidades de tensor nativas, conservando prácticamente el
        $100\%$ de la precisión original.

    \item \textbf{INT8 (Enteros de 8 bits):} Mapea los pesos de 32 bits a valores
        enteros discretos dentro del rango $[-128, 127]$. 
\end{itemize}

Pasar un modelo de detección u OCR desde FP32 a INT8 reduce la huella de memoria
del archivo de pesos en un $75\%$ y triplica la velocidad de inferencia (FPS) en
aceleradores como TensorRT, OpenVINO o ONNX Runtime. Dado que la cuantización
transforma operaciones de multiplicación matricial complejas en aritmética de
enteros de bajo consumo, los dispositivos de borde pueden procesar los flujos de
video con una latencia de milisegundos, experimentando una degradación en la
precisión (mAP) casi imperceptible ($\le 1\%$). Esta optimización permite
desplegar modelos de vanguardia (como la familia YOLO) directamente en las
casetas del estacionamiento, garantizando un control de acceso autónomo, robusto
y eficiente.

\subsection{Valores de referencia y umbrales de aceptación del proyecto}

Las métricas definidas en los apartados precedentes carecen de poder de decisión
mientras no se establezca un valor de referencia que permita calificar el
desempeño de un modelo como suficiente o insuficiente. Un umbral desprovisto de
justificación constituye un criterio arbitrario, por lo que en este apartado se
explicitan los valores adoptados y su procedencia.

La justificación combina dos fuentes complementarias. La primera es el rango de
desempeño reportado en la literatura para tareas equivalentes, que sitúa el
umbral dentro de lo que la comunidad científica considera alcanzable. La segunda
es el requisito operativo del sistema: el nivel de error a partir del cual el
control de acceso deja de ser confiable y exige intervención manual. Cuando la
literatura revisada no reporta un rango directamente comparable, el umbral se
sustenta exclusivamente en la segunda fuente y así se declara de forma explícita.

La Tabla \ref{tab:umbrales} resume los umbrales adoptados, el valor de referencia
que los respalda y su procedencia.

\begin{longtable}{
        P{0.25}
        P{0.10}
        P{0.35}
        P{0.20}
    }
    \caption{Umbrales de aceptación y valores de referencia}
    \label{tab:umbrales}\\

    \toprule
    \rowcolor{gray!20}
    \textbf{Indicador} &
    \textbf{Umbral} &
    \textbf{Valor de referencia} &
    \textbf{Procedencia} \\
    \midrule
    \endfirsthead

    \multicolumn{4}{c}{\textit{Continuación de la Tabla
    \ref{tab:umbrales}}} \\
    \toprule
    \rowcolor{gray!20}
    \textbf{Indicador} &
    \textbf{Umbral} &
    \textbf{Valor de referencia} &
    \textbf{Procedencia} \\
    \midrule
    \endhead

    mAP@50 (detección de vehículos) & $\geq 80\%$ &
    Entre $40{,}9$ y $57{,}5$ de mAP@50--95 en COCO para la familia YOLO26. &
    \cite{jocher2026yolo26, sapkota2026yolo26} \\
    \midrule

    F1-score macro (tipo de vehículo) & $\geq 85\%$ &
    Sin rango directamente comparable en las fuentes revisadas. &
    Criterio operativo \\
    \midrule

    F1-score macro (color predominante) & $\geq 85\%$ &
    Entre $83{,}5\%$ y $97{,}8\%$, con $44{,}4\%$ de imágenes nocturnas mal
    clasificadas en condiciones adversas. &
    \cite{lima2024vehiclecolor} \\
    \midrule

    Recall (detección de placas) & $\geq 90\%$ &
    Entre $91{,}3\%$ y $98{,}8\%$ de exactitud de detección en la revisión
    sistemática de sistemas de reconocimiento de placas. &
    \cite{shashirangana2021alpr, anagnostopoulos2008lpr} \\
    \midrule

    Exactitud por carácter (OCR) & $\geq 90\%$ &
    $89{,}6\%$ de exactitud a nivel de carácter sobre datos reales con
    matrículas del contexto local. &
    \cite{auzabanegas2026blpr} \\
    \midrule

    Exactitud de eventos & $\geq 90\%$ &
    Las fuentes revisadas definen las métricas de seguimiento pero no un umbral
    de aceptación. &
    Criterio operativo \\
    \midrule

    Latencia por evento & $\leq 1$ s &
    Sin rango comparable; condicionada por la exigencia de operación en tiempo
    real. &
    Criterio operativo \\*
    \multicolumn{4}{r}{\footnotesize\textit{Fuente: Elaboración propia}} \\*
    \bottomrule
\end{longtable}

\textbf{Detección de vehículos.} La familia YOLO26 reporta en COCO valores de
mAP@50--95 comprendidos entre $40{,}9$ y $57{,}5$ según la escala del modelo,
métrica que resulta más exigente que el mAP@50 evaluado en este proyecto, pues
promedia diez umbrales de intersección sobre unión en lugar de uno. Considerando
además que el dominio de estudio es sustancialmente más acotado que COCO
---una única escena, condiciones de captura estables y una sola familia de
objetos---, el umbral de $80\%$ de mAP@50 se sitúa por encima del desempeño
generalista reportado y es, por tanto, exigente pero alcanzable.

\textbf{Reconocimiento de color.} El rango reportado para reconocimiento de color
vehicular abarca desde $83{,}5\%$ en condiciones adversas hasta $97{,}8\%$ en
condiciones controladas. El umbral de $85\%$ se ubica deliberadamente en el
extremo conservador de ese rango, próxima al límite inferior, porque las
condiciones de captura del parqueadero incluyen operación nocturna, donde la
misma fuente reporta que cerca de la mitad de las imágenes nocturnas resultan
mal clasificadas.

\textbf{Reconocimiento de placas.} El umbral de Recall adoptado ($90\%$) se sitúa
en el extremo inferior del rango de detección reportado ($91{,}3\%$ a $98{,}8\%$),
lo que lo convierte en el criterio menos exigente del conjunto. En cambio, el
umbral de exactitud por carácter ($90\%$) se apoya en un resultado obtenido sobre
datos reales con matrículas del contexto local, que reporta $89{,}6\%$; adoptar
$90\%$ implica exigir un desempeño igual o superior al reportado, lo que
constituye uno de los criterios más demandantes del proyecto.

\textbf{Tipo de vehículo y determinación de eventos.} Las fuentes revisadas
fundamentan la elección de las métricas ---media armónica para el desbalance
entre clases y las métricas CLEAR MOT para el seguimiento--- pero no establecen
un valor de aceptación universal. El umbral de $85\%$ para el tipo de vehículo y
el de $90\%$ para los eventos se justifican, en consecuencia, por vía operativa:
un error superior a esos niveles implicaría que una proporción apreciable de
vehículos requeriría verificación manual en la barrera, anulando el beneficio de
la automatización.

\textbf{Latencia.} La restricción de un segundo por evento no proviene de la
literatura sino de la operación del recinto: por encima de ese valor, el tiempo
de espera del conductor en la barrera deja de ser aceptable. Su verificación es
experimental y se realiza sobre el hardware de despliegue.

En síntesis, tres de los siete criterios se sustentan simultáneamente en la
literatura y en el requisito operativo, mientras que los cuatro restantes se
sustentan únicamente en el segundo. Esta distinción se declara de forma explícita
para que la evaluación posterior permita distinguir entre un resultado que
alcanza un estándar de la disciplina y uno que satisface una exigencia
operativa del proyecto.
